\chapter{Key Reference}
\label{ch:keyref}

Every key \ccprod{} responds to, arranged by where you are when you press it.

\section{In the worksheet}
\index{keys!worksheet}

This is \msg{READY} or \msg{MODIF} mode: the status line says one of those, and
you are not typing into a cell.

\subsection{Moving}

\begin{cctable}{Key}{Action}{1.6in}{3.22in}
\key{$\uparrow$} & Up one row \\
\key{$\downarrow$} & Down one row \\
\key{$\leftarrow$} & Left one column \\
\key{$\rightarrow$} & Right one column \\
\key{Return} & Down one row \\
\key{Tab} & Right one column \\
\keyc{Shift}{Tab} & Left one column \\
\key{PgUp} & Up one screenful (54 rows) \\
\key{PgDn} & Down one screenful \\
\key{Home} & Column \cellref{A} of this row \\
\keyc{Shift}{Home} & Cell \cellref{A1} \\
\key{End} & Column \cellref{IV} of this row --- \emph{not} the last used
column \\
\end{cctable}

\subsection{Commands}

\begin{cctable}{Key}{Action}{1.6in}{3.22in}
\key{F1} & \menupath{File > New} \\
\key{F2} & \menupath{Edit > Edit cell} \\
\key{F3} & \menupath{File > Open} \\
\key{F4} & \menupath{File > Save} \\
\key{F5} & \menupath{Edit > Undo} \\
\key{F6} & \menupath{File > Save as} \\
\key{F7} & \menupath{File > Import} \\
\key{F8} & \menupath{File > Export} \\
\key{F9} & \menupath{File > Quit} \\
\key{F10} & Open the menu bar \\
\key{F11} & \menupath{Edit > Paste} --- the same as \keyc{Ctrl}{V} \\
\key{F12} & \textbf{Not used.} It belongs to the machine's debugger. \\
\keyc{Ctrl}{C} & \menupath{Edit > Copy} \\
\keyc{Ctrl}{V} & \menupath{Edit > Paste} \\
\keyc{Ctrl}{F} & \menupath{Edit > Find} \\
\key{Backspace} & Clear the current cell \\
\key{Delete} & Clear the current cell \\
\key{Esc} & Nothing \\
Any printable key & Start typing a new entry into the current cell \\
\end{cctable}

\begin{ccnote}
The \menupath{Data}, \menupath{Chart} and \menupath{Sheet} menus have no keys
at all, and neither do \menupath{Insert row} and \menupath{Delete row} on
\menupath{Edit}. Sorting, recalculating, charting, freezing, inserting or
deleting a row, and every worksheet command can only be reached through the
menu bar.
\end{ccnote}

\subsection{Keys that share a code}
\index{keys!duplicate codes}

The Commander X16 sends the same character code for certain control
combinations as it does for named keys, and \ccprod{} cannot tell them apart.
Pressing either of a pair does whatever this manual says the named key does.

\begin{cctablethree}{This}{Is the same as}{Which does}{1.15in}{1.15in}{2.5in}
\keyc{Ctrl}{B} & \key{PgDn} & Down one screenful \\
\keyc{Ctrl}{D} & \key{End} & To column \cellref{IV} \\
\keyc{Ctrl}{I} & \key{Tab} & Right one column \\
\keyc{Ctrl}{M} & \key{Return} & Down one row \\
\keyc{Ctrl}{P} & \key{F9} & Quit \\
\keyc{Ctrl}{Q} & \key{$\downarrow$} & Down one row \\
\keyc{Ctrl}{S} & \key{Home} & To column \cellref{A} \\
\keyc{Ctrl}{T} & \key{Backspace} & Clear the cell \\
\keyc{Ctrl}{U} & \key{F10} & Open the menu bar \\
\keyc{Ctrl}{X} & \keyc{Shift}{Tab} & Left one column \\
\keyc{Ctrl}{Y} & \key{Delete} & Clear the cell \\
\end{cctablethree}

\begin{cccaution}
\keyc{Ctrl}{P} quits the program and \keyc{Ctrl}{T} clears the cell. Neither is
what those combinations do in most software, and the first of them will ask to
discard your changes if you have any.
\end{cccaution}

\section{While typing into a cell}
\index{keys!editing}

This is \msg{EDIT} mode. The editor sees every key first, so nothing in this
table can be overridden by a command.

\begin{cctable}{Key}{Action}{1.6in}{3.22in}
Any printable key & Insert that character at the insertion point \\
\key{$\leftarrow$} & Insertion point left one character \\
\key{$\rightarrow$} & Insertion point right one character \\
\key{Home} & To the start of the line \\
\key{End} & To the end of the line \\
\key{Backspace} & Delete the character to the left \\
\key{Delete} & Delete the character at the insertion point \\
\key{Return} & Store, and move down \\
\key{$\downarrow$} & Store, and move down \\
\key{$\uparrow$} & Store, and move up \\
\key{Tab} & Store, and move right \\
\keyc{Shift}{Tab} & Store, and move left \\
\key{Esc} & Abandon the entry \\
Anything else & Ignored \\
\end{cctable}

The entry is limited to eighty characters. Beyond that, keystrokes are
discarded silently.

\section{With a menu open}
\index{keys!menus}

\begin{cctable}{Key}{Action}{1.6in}{3.22in}
\key{$\leftarrow$} & Previous menu, wrapping round the five \\
\key{$\rightarrow$} & Next menu, wrapping round the five \\
\key{$\uparrow$} & Previous item, wrapping \\
\key{$\downarrow$} & Next item, wrapping \\
\key{Return} & Choose the highlighted item \\
\key{Esc} & Close the menu \\
\key{F10} & Close the menu \\
\key{F1}--\key{F9}, \key{Backspace}, \keyc{Shift}{Home} & Close the menu and do
what that key does in the worksheet \\
\key{PgUp}, \key{PgDn} & Ignored \\
\end{cctable}

\section{In the file list}
\index{keys!dialogues}

\begin{cctable}{Key}{Action}{1.6in}{3.22in}
\key{$\uparrow$} & Previous file. Does not wrap. \\
\key{$\downarrow$} & Next file. Does not wrap. \\
\key{Return} & Open the highlighted file \\
\key{Esc} & Cancel \\
First click on a line & Select it \\
Second click on the same line & Open it \\
Click off the list & Cancel \\
\end{cctable}

\section{In a name box}

The full editor is available: arrows, \key{Home}, \key{End}, \key{Backspace},
\key{Delete} and insertion.

\begin{cctable}{Key}{Action}{1.6in}{3.22in}
\key{Return} & Accept the name \\
\key{Esc} & Cancel \\
\key{Return} on an empty field & Cancel \\
Click on the hint line, left half & The same as \key{Return} \\
Click on the hint line, right half & The same as \key{Esc} \\
\end{cctable}

Names are cut off at twenty-three characters for a file and thirty-one for a
worksheet.

\section{In a message or confirmation box}

\begin{cctable}{Box}{Keys}{1.6in}{3.22in}
Message (\msg{press any key}) & Any key, or any click, dismisses it \\
Confirmation & \key{Y} means yes. \textbf{Any other key means no.} A click on
\msg{[~~Yes~~]} means yes; a click anywhere else means no. \\
\end{cctable}

\section{At the Find prompt}

\begin{cctable}{Key}{Action}{1.6in}{3.22in}
Any printable key & Add a character, up to fifteen \\
\key{Backspace} & Remove the last character \\
\key{Return} & Search \\
\key{Esc} & Abandon \\
Arrow keys & Ignored \\
\end{cctable}

\section{While a chart is on screen}
\index{keys!charts}

\begin{cctable}{Key}{Action}{1.6in}{3.22in}
\key{S} & Save the chart to the card as \fname{CHART.BMP}, then go back to the
worksheet \\
Any other key & Go back to the worksheet \\
\end{cctable}

\key{S} works in either case. See Chapter~\ref{ch:charting}.

\section{At the Delete row question}

\key{Y}, in either case, deletes the row. Any other key cancels.
